\section{Further research: precise next questions}
\label{research:section}

The proved results suggest several exact research problems.
They are listed with the additional information needed for
a solution. None is asserted on numerical evidence alone.

\subsection{Centered harmonic values and global spectral derivatives}

\paragraph{R1. Evaluate mixed even-tail moments in a controlled algebra.}
The invariant-ring theorem includes every product
\[
 (\zeta(2i)-H_n^{(2i)})(\zeta(2j)-H_n^{(2j)})
\]
and its products with odd-order harmonic numbers and powers
of $H_n$. Determine finite formulas for their centered
values at $s=-2r$, beginning with $i=1,j=2$.
The proof of Theorem~\ref{ghp:thm:square} suggests finite
summation by parts followed by Bernoulli coefficient
extraction. The target should specify an ordinary or
multiple-zeta coordinate space at each weight, with an
explicit account of the terms that survive. A useful
strong result would sum an entire family in one convergent
generator, as in \eqref{ghp:eq:polylog-generator}.

\paragraph{R2. Classify shifts other than one half.}
For a fixed positive rational $a\ne1/2$, determine all
polynomials $P$ whose principal parts vanish on the even
nonpositive spectral lattice. Proposition~\ref{ghp:prop:principal}
gives every finite obstruction, but the centered involution
no longer has its simple form. The key question is whether
the infinitely many Bernoulli constraints admit a finite
algebraic description. A classification of simultaneous
conditions at $a$ and $1-a$ may be more natural than a
classification at one noncentral shift.

\paragraph{R3. Separate local regularity from global first derivatives.}
The absence of a principal part is a local asymptotic
statement. The value of the first regular spectral
derivative also contains a global remainder integral.
Find polynomial combinations $P_1,\ldots,P_k$ and exact
coefficients for which these remainder kernels cancel
at a nonpositive even center. A satisfactory answer is
an identity of kernels or a normally convergent sum,
not a numerical relation among already evaluated
derivatives. The centered parity classification narrows
the possible local terms and should reduce this search.

\subsection{Beyond rational cotangent kernels}

\paragraph{R4. Complex powers with moving endpoint exponents.}
Extend the completed generator to
\[
 \operatorname{FP}_x\int_0^1
     \gamma_m^{(p)}(x)(K(x)-z)^{-\lambda}\,dx,
 \qquad z\notin\mathbb R.
\]
The local exponent is $\lambda-p-1$ before logarithmic
differentiation. Thus integer crossings of $\lambda$
create new resonances not present in a fixed repeated
resolvent. First choose a logarithm of $K(x)-z$
continuously along the interval, then prove a joint
meromorphic generator and identify each full resonant
amplitude before expanding either parameter. The rational
integer cases of Corollary~\ref{res:repeated} are exact
normalization data for this problem.

\paragraph{R5. Products of continuous polygamma orders.}
The transform \eqref{res:continuousformula} is linear in
the generalized function $\Psi(v,x)$. For products such
as $\Psi(v,x)\Psi(w,x)$, Fourier convolution introduces
ordered or Tornheim-type coordinates. Determine which
symmetric or reflected combinations reduce to ordinary
polylogarithmic order derivatives and which genuinely
retain a multivariate zeta coordinate. The useful
deliverable is an explicit joint generator together with
the contact terms on every resonant hyperplane.

\paragraph{R6. Unbalanced Barnes functions with fixed normalization.}
The balanced ladder gives combinations of
$\zeta'(-r,x)$ and $\zeta(-r,x)$, and hence normalized
Barnes combinations at the first levels. Seek individual
resolvent transforms of $\log G(x)$ and higher multiple
Gamma functions, with every polynomial-in-$x$ companion
term and normalization constant stated. A functional
equation alone does not determine the function without
the customary asymptotic normalization. One concrete
starting point is to combine \eqref{res:balancedmoment}
with transforms containing $x\log\Gamma(x)$, whose
Fourier coefficients can be computed by one integration
by parts.

\subsection{Entire harmonic interpolation at more than one free order}

\paragraph{R7. A jointly convergent generator for several free orders.}
The canonical gap polynomial is already explicit for
$k$ surviving free indices. Sum arbitrary zero blocks
before, after, and between these indices into a useful
multivariate analytic kernel. The coefficient identity
alone is insufficient: prove a domain in all decoration
variables, the cancellation of simultaneous Lerch
branches, and normal convergence of independent spectral
derivatives. The one-free-order formula
\eqref{harm:zerogenerator} supplies a stringent marginal
specialization for any proposed answer.

\paragraph{R8. The first transverse derivative at a deleted index.}
Differentiate a word in an index that will subsequently
be specialized to $-m$. Its lattice summand acquires
$-x^m\log x$, so a Bernoulli polynomial no longer
represents the gap sum. Identify a finite family of
new gap functions, perhaps in normalized Hurwitz-zeta
derivatives, closed under the same concatenation and
inversion argument. The first target is a fully proved
formula for
\[
 \left.\partial_u\mathcal H_{a,b}(s,u,t)\right|_{u=0}
\]
at arbitrary independent $s,t$. This would extend the
present theorem in a genuinely transverse direction.

\paragraph{R9. Rational endpoints and exact character projections.}
Combine the explicit endpoint coefficient formulas with
finite character decompositions at rational $a,b$.
The expected output is an algorithm expressing the
specified harmonic jets and polynomial Stieltjes moments
in Dirichlet $L$-derivatives and conductor logarithms.
The finite reduction should identify the character
parity, the normalization of the endpoint difference,
and every logarithmic derivative of the conductor
factor. This is particularly suited to exact symbolic
verification.

\subsection{The next Tornheim coordinates and the Gaussian conjectures}

\paragraph{R10. Determine the second-degree symmetric germ.}
The polynomial $J_2$ in \eqref{qtr:J} has two independent
formal coordinates under the named relations.
The diagonal fifth derivative provides only one linear
combination. Obtain one additional independent analytic
identity or a normalized ordinary convergent integral
that determines the other combination. The explicit
deformation following Corollary~\ref{qtr:cyclicimage}
proves why manipulating the existing diagonal and
Euler-plane relations alone cannot finish this task.

\paragraph{R11. Compare additional functional relations with the exact kernel.}
The representation \eqref{qtr:kernelform} is an explicit
target against which another Tornheim relation can be
tested. Expand a proposed functional relation at the
origin, compute its action on the $S$ and $Q$ components,
and state exactly which directions it removes. This
would turn a collection of individual ray identities
into a systematic all-order reduction theory. Any
improved dimension statement must list the added
relations; numerical dependence among sampled
coordinates is not a substitute.

\paragraph{R12. Give an independent Gamma or Barnes representation for \(\Xi\).}
The ordinary integral \eqref{qtr:Xi} is explicit, but it
does not yet explain whether the mixed coordinate
belongs to a smaller familiar constant algebra.
A useful next theorem would derive a different Gamma
or Barnes representation and prove its equality with
the polylogarithmic integral. Such an independent
representation may reveal functional symmetries or
exact cancellations that are hidden in either form.
Simply changing the splitting point in the same
subtracted integral does not provide new information.

\paragraph{R13. Find a functional bridge to the actual \(S_6\) and revised \(S_8\) candidates.}
The short Gaussian reductions remain important open
problems in the source. The completed bounded relation
obstruction shows that a proof must use information
beyond the particular imposed Cayley and extra-row
system. A constructive next step is an explicit
one-parameter identity whose Taylor coefficient is the
target mixed harmonic sum and whose second evaluation
uses the proposed Gaussian polylogarithmic constants.
All boundary terms, algebraic arguments, and logarithm
branches must be fixed before specialization. The
rational and Lerch generators in this article offer
possible tools, but no such bridge to either named
candidate has been proved here.

\subsection{A practical order of attack}

The most immediate exact projects are R1, R7, and R8:
each starts with a finite normal form or a convergent
kernel already supplied here. R10 and R11 are the
natural next steps for the directional-zeta program,
because the obstruction identifies precisely what
new information is missing. R4 and R5 require a more
substantial joint resonance analysis.

For formal verification, the first useful units are
the Bernoulli deletion identities, the canonical gap
normal form, the finite principal-part coefficients,
the quartic coefficient map, and the kernel decomposition
\eqref{qtr:kernelform}. These are exact finite algebraic
statements. Their analytic interfaces should then
be formalized separately: normal convergence,
endpoint subtraction, and the specified continuation
theorems. This separation makes a future machine-checked
development concrete without pretending that the
present numerical diagnostics already supply one.
